\section{Introduction}\label{sec:intro}
\IEEEPARstart{D}{isassembly} separates end-of-life products into components and materials that can re-enter production. Profitable recovery requires more than removing every accessible part. A low-value operation can consume scarce station time, whereas an apparently unprofitable predecessor can provide access to a valuable component. The disassembly line balancing problem (DLBP) therefore connects task selection, precedence, station assignment, and recovery economics. Guo \etal\ review the modeling and optimization of disassembly sequences \cite{Guo2021}. \"Ozceylan \etal\ survey line-balancing formulations and solution approaches \cite{Ozceylan2019}. These perspectives motivate treating selection and allocation as a coordinated decision.

Selective disassembly makes this coordination particularly important. Only part of a product may justify recovery, but required operations and minimum recovery targets still restrict the selected subset. G\"uler \etal\ organize developments in partial disassembly line balancing \cite{Guler2024}. Laili \etal\ compare models and algorithms for multi-objective disassembly balancing \cite{Laili2020}. He \etal\ examine uncertainty and links to supply-chain decisions \cite{He2024}. These reviews show that problem variants differ in both feasible operations and performance criteria. A comparison of solution methods consequently needs a fixed mathematical model and identical input data. The present work uses deterministic task data and one profit objective so that formulation and certification effects remain identifiable.

Uncertainty models are application-specific. Wang \etal\ study stock-market ambiguity through a Markov regime-switching model \cite{Wang2021Markov}. That financial setting differs from disassembly failures and processing-time uncertainty; the present formulation uses fixed task data.

Human--robot collaboration adds another layer of operational choice. Human execution can offer flexibility, robot execution can provide a different time--cost combination, and joint execution uses both resources. Xu \etal\ study collaborative line balancing with a safe strategy \cite{Xu2021}. Parsa and Saadat investigate collaborative planning for end-of-life product disassembly \cite{Parsa2021}. Xiao and Huang review collaboration in uncertain and dynamic disassembly \cite{Xiao2024}. Yin \etal\ formulate multi-man--robot line balancing with mathematical programming and a genetic Jaya algorithm \cite{Yin2025}. Qin \etal\ examine multiple-product hybrid disassembly line balancing with human--robot collaboration \cite{Qin2025Hybrid}. These studies support explicit mode decisions, while their safety, workforce, and uncertainty assumptions need not coincide with those of a deterministic shared-resource model.

Layout also changes which decisions interact. Wu \etal\ consider two-sided balancing with human--robot interaction constraints \cite{Wu2023}, and Chu and Chen study partial disassembly in a two-sided cooperative setting \cite{Chu2025}. Al-fandi and Almasarwah examine mixed-model two-sided disassembly \cite{Alfandi2025}. Huang \etal\ introduce activation decisions in a collaborative U-shaped line \cite{Huang2025}. Fang \etal\ address position-constrained collaborative planning \cite{Fang2024}. Batta\"ia and Dolgui review hybrid line-balancing structures \cite{Battaia2022}, while Boysen \etal\ review broader assembly-balancing developments \cite{Boysen2022}. Wang \etal\ apply reinforcement learning to hybrid disassembly line balancing \cite{Wang2024RL}. These contributions indicate why layout assumptions require explicit treatment. Here a physical circular conveyor has a fixed entry and an ordered, single-pass task sequence. Its geometric shape does not introduce another product lap or a cyclic precedence decision.

Profit-oriented collaboration also interacts with operation-specific restrictions. Guo \etal\ study hybrid partial destructive disassembly with noise and component accessibility considerations \cite{Guo2024}. Zhu \etal\ consider remanufacturing demand and hazardous components \cite{Zhu2024}. Goksoy Kalaycilar \etal\ address hazardous task failures \cite{GoksoyKalaycilar2024}. Zhang \etal\ maximize collaborative disassembly profit with an improved grey wolf optimizer \cite{Zhang2025}. Wu \etal\ investigate stochastic processing and carbon-emission objectives \cite{Wu2024}, and Guo \etal\ consider stochastic operation times \cite{Guo2024a}. Xu \etal\ develop a multi-objective discrete brainstorming optimizer for stochastic multiple-product robotic disassembly balancing with failures \cite{Xu2023Failure}. The present model retains task-level recovery value and mode-dependent execution cost. It does not import hazard, failure, carbon, or stochastic constraints from these distinct variants.

Human workload requires a similarly precise definition. Guo \etal\ consider a human fatigue index in U-shaped disassembly balancing \cite{Guo2023}. Otto and Scholl incorporate ergonomic risks into assembly balancing \cite{Otto2011}. Xu and Han explicitly include worker rest time in a two-sided disassembly model \cite{Xu2024}. Zhou \etal\ consider worker fatigue in multi-factory remanufacturing \cite{Zhou2025}. This work limits the cumulative fatigue load at each human-equipped station over one representative cycle. It includes no recovery period and makes no claim about fatigue across an entire shift. This bounded interpretation keeps the ergonomic assumption consistent with the implemented additive-load constraints.

Multiple lines create a further coordination problem even when every individual plan is feasible. Their combined human and robot deployments can exceed the available pools. Guo \etal\ study resource sharing in multiproduct human--robot collaborative hybrid lines \cite{Guo2025}. Gu \etal\ optimize multi-factory remanufacturing with shared transportation \cite{Gu2025}. The shared quantities here are fixed station-level human and robot deployments, rather than transport movements. All task, precedence, recovery, time, and fatigue requirements remain within a line. Only two aggregate resource rows connect different lines. This structure suggests a decomposition that preserves complete local plans.

A compact formulation in mixed-integer linear programming (MILP) provides a direct baseline and an independent definition of feasibility. Its linear programming (LP) relaxation, however, can combine fractional task assignments inside stations. A whole-line formulation instead expresses each line as a convex combination of complete feasible integer plans. Dantzig and Wolfe establish the underlying decomposition principle \cite{Dantzig1960}, and L\"ubbecke and Desrosiers explain pricing and dual bounds in column generation (CG) \cite{Lubbecke2005}. Local integer convexification can strengthen the relaxation, but the corresponding line-pricing problem is itself a MILP. The computational benefit therefore depends on the balance between local pricing difficulty and global coordination, rather than total task count alone.

Integer recovery presents a separate issue. An optimal restricted master LP neither guarantees an integer combination of its columns nor proves the original integer optimum. Barnhart \etal\ describe branch-and-price as a framework for integrating pricing into integer search \cite{Barnhart1998}. Pessoa \etal\ demonstrate generic exact-solver development for routing-related problems \cite{Pessoa2020}. These works establish useful methodological distinctions, rather than a claim that the present implementation performs a branch-and-price tree search. For the two shared resource counts considered here, resource-profile value functions permit an alternative integer coordination mechanism using dynamic programming.

Algorithmic improvements also require evidence about where time and proof gaps arise. du Merle \etal\ study stabilized CG \cite{duMerle1999}, and Maher and R\"onnberg investigate pricing schemes that promote high-quality integer solutions \cite{Maher2023}. Their methods illustrate that pricing behavior and integer solution quality are distinct design concerns. This work uses exact line pricing, verified restricted integer recovery, and demand-driven queries of unresolved resource profiles. It reports valid upper bounds alongside feasible profits and distinguishes LP convergence from integer bound closure. A retained difficult case with few large lines identifies a regime where decomposition loses its advantage.

The choice of a column also affects practical feasibility. With station columns, individually legal configurations can omit required tasks, duplicate a task, or disagree on cross-station precedence when combined. A whole-line column resolves these interactions before entering the master. The remaining resource conflict can be described by two integer counts, without replacing any original task constraint by a penalty. This design trades a harder local solve for a smaller coordination problem. It motivates reporting both pricing effort and integer certification, and testing configurations with different line-block sizes. It also makes the limits of the approach visible: a compact model can remain preferable when local blocks become too difficult or shared resources impose little coordination.

This work makes the following new contributions:
\begin{enumerate}
\item A whole-line reformulation is developed for selective multi-line collaborative disassembly with cumulative station fatigue. Complete line plans preserve all local compatibility requirements, and bidirectional mappings establish integer equivalence with the compact model. The relaxation hierarchy explains the source and limits of formulation strength.
\item A solution framework combines two-phase CG, independently checked integer recovery, and demand-driven resource-profile bounds. Explicit pricing states and repaired dual bounds distinguish feasible solutions, complete LP proofs, and original integer optimality. Profile coordination can close an integer gap without asserting LP convergence.
\item A reproducible comparison uses frozen paired synthetic inputs and reports profits, valid bounds, proof states, pricing rounds, and execution times. Integer certification is faster on the tested four- and six-line configurations. Resource-ample and large-block cases delimit this empirical finding.
\end{enumerate}

Section~\ref{sec:model} defines the problem and compact model. Section~\ref{sec:method} presents the proposed algorithm and its bounds. Section~\ref{sec:experiments} reports the computational evidence. Section~\ref{sec:conclusion} states the conclusions and remaining questions.
